\documentclass[11pt,a4paper]{article}

\usepackage{geometry}
\geometry{margin=1in}

% 1. Bilingual font configuration
\usepackage{greekenglish}

% 2. Dark code syntax highlighting
\usepackage{code}

\title{\textbf{Research Tools: LaTeX Libraries Demonstration}}
\author{Asterinos}

\begin{document}

\maketitle

\section{Introduction / Εισαγωγή}
This document demonstrates the two custom packages in the \texttt{latex/} folder:
\begin{itemize}
    \item \texttt{greekenglish.sty}: Enables typing both English and Greek seamlessly using Unicode fonts with LuaLaTeX / XeLaTeX.
    \item \texttt{code.sty}: Provides stylish dark-mode syntax highlighted code environments for various languages.
\end{itemize}

\subsection{Bilingual Text Switching (Άμεση εναλλαγή γλωσσών)}
Here is an example use of English and Greek text:
\begin{quote}
    The quick brown fox jumps over the lazy dog. 
    Η γρήγορη καφέ αλεπού πηδά πάνω από το τεμπέλικο σκυλί. 
    Η επιστήμη των υπολογιστών και οι σύγχρονες τεχνολογίες επεξεργασίας δεδομένων.
    Computer science and modern data engineering technologies.
\end{quote}
No commands were required.

\section{Code Highlighting Examples}

\subsection{Python Environment}
\begin{python}
# Python sample code
def fibonacci(n: int) -> list[int]:
    """Generate Fibonacci sequence up to n numbers."""
    if n <= 0:
        return []
    seq = [0, 1]
    while len(seq) < n:
        seq.append(seq[-1] + seq[-2])
    return seq[:n]

if __name__ == "__main__":
    print(f"Fibonacci(8): {fibonacci(8)}")
\end{python}

\subsection{SQL Environment}
\begin{sql}
-- Query user research analytics
SELECT 
    u.id, 
    u.username, 
    COUNT(r.id) AS total_runs,
    AVG(r.execution_time_ms) AS avg_time_ms
FROM users u
LEFT JOIN research_runs r ON u.id = r.user_id
WHERE r.created_at >= '2026-01-01'
GROUP BY u.id, u.username
ORDER BY total_runs DESC;
\end{sql}

\subsection{JavaScript Environment}
\begin{javascript}
// Modern async fetch function
async function fetchResearchData(endpoint) {
    try {
        const response = await fetch(`https://api.research.local/${endpoint}`);
        if (!response.ok) {
            throw new Error(`HTTP error! status: ${response.status}`);
        }
        const data = await response.json();
        return data;
    } catch (err) {
        console.error("Failed to fetch research metrics:", err);
    }
}
\end{javascript}

\end{document}
